% !PRINTED 95
\weberpage{95}

If two columns in $A$ or in $B$ are interchanged, the elements $c_i^{(h)}$ do not change, but only interchange among themselves. If, however, two rows in $A$ or in $B$ are interchanged, then the $c_i^{(h)}$ change because the factors in the individual products of the sum are grouped differently; but if \gesperrt{corresponding} interchanges are made simultaneously in the rows of $A$ and of $B$, the $c_i^{(h)}$ remain unchanged, since only the individual terms of the sum are thereby interchanged.

By turning the rows into columns in $A$ or in $B$, or in both at once, one obtains three further different ways of forming the product of two determinants in determinant form. This can also be expressed as follows: \gesperrt{To form the product of two determinants, one may multiply the elements of the separate rows or columns of one factor by the corresponding elements of the rows or columns of the other factor, add the products, and regard these sums of products as the elements of a new determinant.}

The multiplication rule can be applied to a product of two determinants with different numbers of elements by transforming the determinant with the smaller number of rows, by the theorem § 22, IX, into another with more rows.

\begin{center}
{\large \S.\ 28.}

\vspace{0.5em}
{\large\gesperrt{Determinants of Subdeterminants.}}
\end{center}

We immediately make an application here of the multiplication law of determinants. Let, as before,
\[
A=\begin{vmatrix}
 a_1^{(1)}&a_2^{(1)}&\cdots&a_n^{(1)}\\
 a_1^{(2)}&a_2^{(2)}&\cdots&a_n^{(2)}\\
 \cdots&\cdots&&\cdots\\
 a_1^{(n)}&a_2^{(n)}&\cdots&a_n^{(n)}
\end{vmatrix},\tag{1}
\]
and let
%CONT>>
